\chapter{Ingegnerizzazione del Sistema, REST API e Containerizzazione}

\section{Backend REST API (FastAPI)}
Il backend del sistema è realizzato in Python mediante il framework **FastAPI**. Il servizio interagisce con MongoDB Gold ed espone i seguenti endpoint principali:
\begin{itemize}
    \item \texttt{GET /cases}: Elenco dei casi clinici registrati.
    \item \texttt{GET /cases/\{case\_id\}/series}: Serie temporale del caso selezionato, con supporto all'aggregazione temporale lato database (\textit{downsampling}).
    \item \texttt{POST /cases/\{case\_id\}/detect}: Esecuzione al volo dei modelli di anomaly detection e salvataggio dei risultati.
\end{itemize}

\section{Dashboard Web Interattiva}
Per la visualizzazione intuitiva dei dati è stata realizzata una dashboard web frontend in HTML5/CSS3/JavaScript (con libreria Chart.js).
La dashboard permette all'operatore sanitario o all'analista di:
\begin{itemize}
    \item Selezionare interattivamente un caso dal registro.
    \item Esplorare i grafici temporali sovrapposti delle cinque metriche vitali.
    \item Avviare l'anomaly detection e visualizzare visivamente i marcatori delle anomalie sovrapposti ai tracciati.
\end{itemize}

\section{Containerizzazione ed Orchestrazione con Docker}
L'intero ecosistema è completamente containerizzato ed orchestrato tramite **Docker** e **Docker Compose**:
\begin{itemize}
    \item \textbf{Container \texttt{vitaldb\_mongo}}: Istanza di MongoDB 8.0 avviata automaticamente come Replica Set (\texttt{--replSet rs0}) con verifica di salute nativa (\textit{healthcheck}).
    \item \textbf{Container \texttt{vitaldb\_api}}: Istanza dell'API FastAPI basata su Python 3.12-slim, isolata nella rete interna di Docker.
\end{itemize}

L'avvio dell'intero ambiente viene effettuato con un singolo comando:
$$\texttt{docker compose up --build -d}$$

\section{Conclusioni e Sviluppi Futuri}
Il progetto ha dimostrato con successo l'efficacia dell'integrazione tra moderne pipeline di Data Engineering ed algoritmi avanzati di Anomaly Detection per il dominio sanitario. 

Tra i possibili sviluppi futuri figurano l'integrazione di motori di streaming in tempo reale (ad es. Apache Kafka) e l'estensione del tracciamento ad ulteriori indicatori fisiologici.
